\section*{Задача 4. Гномьи сундуки (Средняя)}
\textbf{Условие.} В сокровищнице короля гномов хранится $N$ сундуков, занумерованных от $1$ до $N$. В каждом сундуке лежит некоторая неотрицательная целая сумма денег. Гномы используют собственную систему счисления, которая отличается от привычной десятичной двумя ключевыми особенностями:
\begin{enumerate}
    \item \textbf{Обратный порядок разрядов.} Цифры записываются начиная с младшего разряда. Если число в обычной записи имеет вид $\overline{d_k \dots d_1 d_0}$, то у гномов оно записывается как строка символов, соответствующих $d_0, d_1, \dots, d_k$. Значение числа вычисляется по формуле:
    $$ S = \sum_{i=0}^{k} \text{val}(d_i) \cdot 10^i, $$
    где $\text{val}(d_i)$ -- числовое значение $i$-й цифры.
    \item \textbf{Набор цифр.} Гномы не используют цифру $0$ для записи ненулевых чисел. Вместо неё используется цифра $10$. Цифра $0$ применяется исключительно для обозначения самого числа ноль, которое записывается как \texttt{()}. Для обозначения цифр от $1$ до $10$ используются следующие последовательности символов:
\end{enumerate}
\begin{center}
\begin{tabular}{c|c|c|c|c|c|c|c|c|c|c}
$1$ & $2$ & $3$ & $4$ & $5$ & $6$ & $7$ & $8$ & $9$ & $10$ \\ \hline
\texttt{>} & \texttt{>!} & \texttt{>!!} & \texttt{>!!!} & \texttt{>?} & \texttt{<} & \texttt{<!} & \texttt{<!!} & \texttt{<!!!} & \texttt{<?}
\end{tabular}
\end{center}
Символы в записи числа идут подряд без пробелов и разделителей. Поскольку каждая цифра кодируется строкой разной длины (от $1$ до $4$ символов), для перевода гномьей записи в число необходимо аккуратно выделять разряды, двигаясь от начала строки (младший разряд) к концу (старший разряд).

\textbf{Задание.} Вам даны суммы, лежащие в сундуках с номерами $1, 2, \dots, N$, записанные в гномьей системе. Требуется найти номера двух сундуков $K$ и $L$ ($1 \le K < L \le N$), такие что абсолютная разность лежащих в них сумм $|S_K - S_L|$ максимально возможна. Если существует несколько пар с одинаковой максимальной разностью, выберите ту, для которой сумма номеров $K + L$ минимальна. Найденные номера $K$ и $L$ необходимо вывести в гномьей системе счисления.

\textbf{Примечания.}
\begin{itemize}
    \item Длина записи суммы может достигать $100$ символов, поэтому числа могут быть очень большими. Стандартные целочисленные типы данных не подойдут; рекомендуется использовать арифметику больших чисел или лексикографическое сравнение нормализованных строк.
    \item Для максимизации $|S_K - S_L|$ достаточно найти сундук с максимальной суммой и сундук с минимальной суммой среди всех $N$ сундуков.
    \item Гарантируется, что $N \ge 2$, поэтому искомая пара всегда существует.
\end{itemize}

\textbf{Формат ввода.} В первой строке записано целое число $N$ ($2 \le N \le 500$) -- количество сундуков. В следующих $N$ строках содержатся записи сумм в гномьей системе (длина каждой строки не превышает $100$ символов). Строки идут в порядке номеров сундуков: первая строка после $N$ соответствует сундуку $1$, вторая -- сундуку $2$, и т.д.

\textbf{Формат вывода.} В первой строке выведите гномью запись номера $K$, во второй строке -- гномью запись номера $L$.

\textbf{Пример.}
\begin{examplebox}
  \begin{minipage}[t]{0.48\linewidth}
    \textbf{Ввод:}\\
    \begin{verbatim}
2
()
>!!!>!<?>
    \end{verbatim}
  \end{minipage}%
  \hfill%
  \begin{minipage}[t]{0.48\linewidth}
    \textbf{Вывод:}\\
    \begin{verbatim}
>
>!
    \end{verbatim}
  \end{minipage}
\end{examplebox}